% !TeX spellcheck = pt_BR

%%%%%%
%
%% Criado por Lucas dos Santos Costa
%% Instituto Nacional de Telecomunicações - Inatel (https://inatel.br/home/)
%% Centro de Referência em Radiocomunicações - CRR (https://www.inatel.br/crr/)

%% Finalizado em julho de 2019.
%% Última atualização: 23/08/2019.
%% Última atualização: 15/07/2020.
%% Última atualização: 25/05/2021.
%% Última atualização: 10/02/2022.
%% Última atualização: 22/11/2023.

%%% A utilização deste Template é bastante simples e algumas soluções para possíveis problemas podem ser encontradas no próprio documento, códigos, e, principalmente, no ``Curso de Introdução ao LaTeX com o TeXstudio'' disponibilizado na pasta ``Curso LaTeX''. 

%-> Sempre use a pasta IEEEtran para tirar dúvidas sobre os padrões utilizados nas referências bibliográficas, equações, tabelas, artigos, etc..

%-> O comando \lipsum[1-2] apenas inseri parágrafos contendo textos aleatórios.
%   Portanto, onde encontrá-lo, comente a respectiva linha para inserir seu texto.

%-> Todos os pacotes (configurações) utilizados no Template estão na pasta 'Preambulo'.

%-> Quando precisar encontrar referências de artigos ou livros no formato bibtex, uma boa opção é usar o 'Google Scholar (https://scholar.google.com.br/)'.
    %Para isso, digite o nome do artigo, por exemplo, no campo de pesquisas e, encontrando-o, selecione a opção aspas duplas ("), escolha Bibtex e copie e cole a referência no formato bibtex para o seu banco (myreferences.bib).
    %Sempre verifique se está tudo dentro do padrão IEEE!
    %Se o artigo (por exemplo) estiver na plataforma do IEEE a referência no formato bibtex pode ser encontrada na própria página do artigo.
    %Para artigos de outras revistas verifique a disponibilidade na respectiva página.

%-> Na pasta 'Curso LaTeX' está disponível um arquivo PDF contendo um curso básico de introdução ao LaTeX.
    %Este material possui informações que podem ser importantes para a confecção de sua dissertação de mestrado/tese de doutorado.
    
%-> Os itens das Listas de Abreviaturas, Siglas, Símbolos e Palavras/termos técnicos (Glossário) se encontram na pasta 'AbrsSigsSymsGlos'.

%-> O TEMPLATE FOI DESENVOLVIDO PARA QUE SEJA FEITA A IMPRESSÃO DO TIPO FRENTE E VERSO.
%POR ISSO FOI CONSIDERADO O PADRÃO USADO EM LIVROS PARA RECUO DE MARGENS EM PÁGINAS PARES E ÍMPARES E A INICIALIZAÇÃO DE CAPÍTULOS EM PÁGINAS ÍMPARES, O QUE ADICIONA UMA PÁGINA EM BRANCO AO FINAL DOS CAPÍTULOS QUE FINALIZAM EM PÁGINAS ÍMPARES!
%Para desfazer uma dessas configurações, retire seu respectivo comando, 'twoside' ou 'openright', de \documentclass[12pt,a4paper,twoside,openright]{report}.

%-> Para definir o espaçamento entre linhas como i) Simples, ii) Um e Meio e iii) Duplo, use, respectivamente, os comandos i) \singlespacing, ii) \onehalfspacing e iii) \doublespacing.
%OBS1: 
    %Estes comandos aplicam o respectivo espaçamento em todo o texto que o sucede! 
    %Nas próximas linhas, por exemplo, o comando \onehalfspacing foi utilizado após o comando \pagenumbering{arabic} para definir o espaçamento entre linhas como 'Um e Meio' em todo o texto entre o Capítulo 1 e os Apêndices.
    %Após os Apêndices, no entanto, o comando \singlespacing foi utilizado antes do comando \include{Glossario/Glossario} para definir o espaçamento entre linhas usado no Glossário como 'Simples'.
%OBS2:
    %Este Template utiliza o Espaçamento Simples por padrão.
%%%%%%


%%%%%%%%%%%%%%%%%%%%%%%%%%%%%%% preâmbulo %%%%%%%%%%%%%%%%%%%%%%%%%%%%%%%
\documentclass[12pt,a4paper]{report}

%%%
%\includeonly{Capitulos/Cap1,Capitulos/Cap2}%Descomente esta linha para compilar apenas o(s) capítulo(s) de interesse
%%%

\input{Preambulo/Configuracoes}
\input{AbrsSigsSymsGlos/AbrsSigsSymsGlos}

%%%%%%%%%%%%%%%%%%%%%%%%%%%%%%% pré-texto %%%%%%%%%%%%%%%%%%%%%%%%%%%%%%%
\begin{document}
	%%%
	\pagenumbering{alph}
	%%%
	
	%%% Capa
	    \includepdf{Capa/CapaTex.pdf}\null\newpage%comente esta linha, se preferir, quanto estiver usando o comando \includeonly{filelist}
	
	%%% ATENÇÃO!!!
    \includepdf[pages={1-}]{Infos/LEIA_ME_PRIMEIRO}%comente esta linha após ler as informações contidas no documento 'LEIA_ME_PRIMEIRO'!
    %%%
    
	%%% Folha de Rosto
		\include{FolhaRosto/FolhaRosto}
	%% Ficha Catalográfica (A ficha catalográfica é fornecida pela biblioteca após a defesa do trabalho)
		\includepdf{FichaCat/FichaCat}%comente esta linha, se preferir, quanto estiver usando o comando \includeonly{filelist}
		
	%%%
	\pagenumbering{roman}
	%%%
	
	%%% Folha de Aprovação
		\include{FolhaAprovacao/FolhaAprovacao}
		%\includepdf{FolhaAprovacao/FolhaAprovacao.pdf}\null\newpage%Depois da defesa, se preferir troque o comando anterior por este para incluir a folha de aprovação assinada pela banca (lembre-se de fazer o upload do arquivo 'FolhaAprovacao.pdf' para a pasta 'FolhaAprovacao'!).
		%Fito esto, lembre-se também de conferir toda a paginação do documento para verificar se não houve alterações/erros!
	%%% Epígrafe
		\include{Epigrafe/Epigrafe}
	%%% Dedicatória
		\include{Dedicatoria/Dedicatoria}
	%%% Agradecimentos
		\include{Agradecimentos/Agradecimentos}
	%%% Sumário	
		\include{Sumario/Sumario}
	%%% Lista de Figuras	
		\include{ListaFiguras/ListaFiguras}
	%%% Lista de Tabelas	
		\include{ListaTabelas/ListaTabelas}
	%%% Lista de Abreviaturas e Siglas
		\include{ListaAbrevSiglas/ListaAbrevSiglas}
	%%% Lista de Símbolos
		\include{ListaSimbolos/ListaSimbolos}
	%%% Resumo
		\include{ResumoAbstract/Resumo}	
	%%% Abstract
	    \acresetall%zera o contador de citação dos acrônimos
		\include{ResumoAbstract/Abstract}
		\cleardoublepage
	
	%%%
	\pagenumbering{arabic}
	\onehalfspacing
	%%%
	
	%%%%%%%%%%%%%%%%%%%%%%%%%%%%%%% texto %%%%%%%%%%%%%%%%%%%%%%%%%%%%%%%%%%%
	\acresetall%zera o contador de citação dos acrônimos
	%%% Introdução
		\include{Capitulos/Cap1}
	%%% Corpo do Trabalho
		\include{Capitulos/Cap2}
		\include{Capitulos/Cap3}
	%%% Conclusões
	    \acresetall%zera o contador de citação dos acrônimos
		\include{Capitulos/Cap4}
	%%%%%%%%%%%%%%%%%%%%%%%%%%%%%%% pós-texto %%%%%%%%%%%%%%%%%%%%%%%%%%%%%%%
	%%% Anexos
	    \include{AnexosApendices/Anexos}
	%%% Referências Bibliográficas
		\include{Referencias}
	%%% Apêndices
		\include{AnexosApendices/Apendices}
	%%% Glossário
	    \singlespacing
		\include{Glossario/Glossario}
\end{document}